\section{Introduction}
\label{sec:intro}

Weight-space model merging has emerged as a foundational technique for constructing unified, multi-task neural networks from a set of specialized, fine-tuned expert models without additional retraining costs \cite{Wortsman2022, Ilharco2022, Yadav2023, Liang2023}. While static merging techniques (e.g., Task Arithmetic \cite{Ilharco2022} or TIES-Merging \cite{Yadav2023}) find a single, consensus set of weights, recent work has pivoted toward \emph{dynamic model merging} \cite{PredecessorT4S6, PredecessorT4S10}. In this paradigm, routing networks process incoming samples at test time and dynamically blend expert parameters on the fly based on representation activations.

Despite their intuitive appeal, we argue that the current state of dynamic routing in model merging has succumbed to a trend of hyper-parameterization and needlessly convoluted architectural metaphors. We identify two severe, overlooked failure modes that plague these designs:

\begin{figure*}[t]
  \centering
  \begin{subfigure}[b]{0.48\textwidth}
    \centering
    \includegraphics[width=\textwidth]{regularization_impact.png}
    \caption{Impact of $L_2$ regularization on unregularized baseline routing methods under homogeneous streams.}
    \label{fig:reg_impact}
  \end{subfigure}
  \hfill
  \begin{subfigure}[b]{0.48\textwidth}
    \centering
    \includegraphics[width=\textwidth]{l3_comparison.png}
    \caption{Comparison of unconstrained layer-wise routing ($L^3$-Router) vs. a global unregularized Linear Router.}
    \label{fig:l3_comparison}
  \end{subfigure}
  \caption{Empirical deconstruction of over-parameterized dynamic routers on our synthetic sandbox. Left: Simple classical $L_2$ regularization eliminates the catastrophic OOD (SVHN) collapse of wave-inspired QWS-Merge, matching or exceeding its performance with standard linear primitives. Right: Demonstrating Layer-Averaging Collapse; unregularized multi-layer routers are systematically outperformed by a simple global, single-layer Linear Router, highlighting parameter redundancy.}
  \label{fig:deconstruction}
\end{figure*}

\textbf{1. Optimization Bloat and Out-of-Distribution Collapse:} Sophisticated routing schemes like Quantum Wavefunction Superposition Merging (QWS-Merge) \cite{PredecessorT4S10} employ complex wave phase activations and multi-layer structures trained via iterative optimization (AdamW \cite{Loshchilov2017}). We demonstrate that these "quantum" metaphors are a design illusion; they suffer from transductive overfitting \cite{PredecessorT2S1} and collapse catastrophically when presented with out-of-distribution (OOD) tasks. For example, QWS-Merge drops to $10.00\%$ accuracy on SVHN (see \cref{fig:reg_impact}). Standard classical $L_2$ regularization applied to a basic linear layer easily replicates or outperforms these convoluted designs. Furthermore, we expose a phenomenon we term \textbf{Layer-Averaging Collapse}: because model merging is typically applied to a single classification head, layer-wise routing networks collapse to a redundant single-layer search space, meaning that a simple global, single-layer Linear Router systematically outperforms layer-wise alternatives (\cref{fig:l3_comparison}).

\textbf{2. Heterogeneity Collapse in Mixed-Task Streams:} Under realistic deployment scenarios, input streams are rarely homogeneous. When a batch contains a mixture of different tasks, dynamic routers must compute sample-wise coefficients and average them across the batch dimension to comply with hardware accelerator requirements. This averaging forces the coefficients toward a flat, uniform distribution, resulting in \textbf{heterogeneity collapse} where the benefits of dynamic routing are completely lost (e.g., the Linear Router drops from $51.00\%$ to $43.40\%$). Some existing models seem robust to this, but we show this is a \textbf{Robustness-Accuracy Illusion}---their simplex-normalization (Softmax) constraints simply force routing coefficients to a flat, mediocre average in \emph{all} scenarios, sacrificing high homogeneous accuracy for artificial "robustness".

In this paper, we adopt the strict philosophy of \textbf{The Minimalist} (Occam's razor). Rather than engineering increasingly complex, heavily-parameterized multi-layer routing networks to survive mixed-task streams, we resolve the issue by (a) stripping away $100\%$ of the routing parameters, and (b) handling stream heterogeneity at the data level. Crucially, our framework operates under the Parameter-Efficient Fine-Tuning (PEFT/LoRA) paradigm to ensure spatial VRAM viability, keeping all active expert parameters as low-rank adapters ($<1\%$ footprint) and capping model-serving memory overhead at a strict 1.04$\times$ base model size. We explicitly highlight that this spatial dependency on PEFT is a non-negotiable, strict prerequisite; full-parameter dynamic weight merging on-the-fly would require keeping all $K$ expert networks fully loaded in GPU memory concurrently (over 70 GB of VRAM for $K=4$ LLaMA-7B experts) or suffering from massive PCIe host-to-device transfer latency (over 5,000 ms per forward pass), rendering dynamic merging non-viable for real-time streaming. In contrast, lightweight LoRA adapters allow kernel-fused dynamic parameter updates in under a millisecond.

We introduce \textbf{Parameter-Free Subspace Routing (PFSR)} coupled with \textbf{Micro-Batch Homogenization (MBH)}. First, PFSR projects high-dimensional intermediate features from the penultimate layer onto a low-dimensional task coordinate subspace by computing maximum cosine similarity against the frozen, pre-trained expert classification weights. It applies a temperature-scaled Softmax directly to these projected task coordinates to generate merging coefficients, requiring \textbf{zero trainable parameters}, \textbf{zero training epochs}, and \textbf{zero calibration data} in standard expert registries (with only a minor, low-resource calibration split reintroduced when fitting unsupervised non-classification centroids or density-based OOD mixture models). Second, MBH resolves heterogeneity collapse at the stream level by dynamically partitioning mixed-task batches into homogeneous micro-batches on the fly, performing specialized inference on each, and re-sorting the final outputs. 

We make the following contributions:
\begin{itemize}
    \item \textbf{Empirical Deconstruction of QWS-Merge:} We prove that wave-inspired routing metaphors are highly unstable and collapse on OOD tasks, and show that classical $L_2$ regularization on a global linear router delivers superior stability and accuracy.
    \item \textbf{Analytical and Empirical Analysis of Layer-Averaging Collapse:} We demonstrate that layer-specific routing coefficients collapse under classification head constraints, rendering multi-layer routing architectures redundant.
    \item \textbf{Introduction of PFSR + MBH:} We present a completely non-parametric, zero-shot merging framework. Our method achieves a high \textbf{75.00\% Joint Mean accuracy} under homogeneous streams, and a robust \textbf{71.60\% Joint Mean accuracy} under heterogeneous mixed-task streams, with absolutely zero parameter overhead.
    \item \textbf{Diverse Real-World Evaluations:} We validate our framework on standard computer vision (Vision Transformers on DomainNet) and large-scale NLP experts (LLaMA-7B on Math, Coding, Translation, and Instruction-Following), demonstrating exceptional scalability and recovering up to $97.5\%$ of the expert standalone ceiling with zero parameter overhead.
    \item \textbf{Integration with PEFT (LoRA) for VRAM Viability:} We ground our dynamic merging framework under the Parameter-Efficient Fine-Tuning (PEFT) paradigm, demonstrating that keeping expert weights as lightweight LoRA adapters ensures a strict $1.04\times$ memory footprint and resolves hardware memory constraints.
    \item \textbf{Instantaneous Dynamic Task Adaptation:} Because PFSR is completely training-free and zero-shot, our framework allows practitioners to register or retire task experts on the fly in massive model hubs with absolutely zero retraining or joint calibration, proving highly valuable for dynamic production registries.
\end{itemize}
